%% =====================================================================
%% 江苏大学本科毕业设计（论文）2026 版 LaTeX 模板
%%
%% 依据：《江苏大学毕业设计（论文）写作规范（2026）》（附件 2）
%%       《江苏大学毕业设计（论文）通用格式模板（理工农医教）2026 版》（附件 3）
%%
%% 编译：xelatex main → bibtex main → xelatex main → xelatex main
%%       目录与交叉引用需要编译两遍。
%%
%% 本文件各页顺序与附件 3 完全一致：
%%   封面 → 原创性声明 → 中文摘要 → 英文摘要 → 目录 → 正文 → 结论
%%        → 参考文献 → 致谢 → 附录
%% =====================================================================
\documentclass{UJSThesis2026}

\begin{document}

%% ---------- 封面 ----------
% !TEX root = ../main.tex
% =====================================================================
% 封面
%
% 版式与文字均对齐附件 3 封面：
%   顶部 校徽 + “江苏大学”手写体 + “JIANGSU UNIVERSITY”（带下划线）+
%        “本科毕业设计（论文）”整块做成**一张图片** figures/ujudt-header.png，
%        由附件 3 的 PDF 第 1 页原样裁出（见 UJSThesis2026.cls 中的说明）。
%   中文题目 二号黑体居中，可分两行
%   英文题目 16pt Times New Roman 加粗居中
%   学院名称／专业班级／学生姓名／学生学号／指导教师姓名／指导教师职称：
%     标签四号宋体，宽 93.9pt 两端对齐，六个“：”竖直对齐；
%     取值四号宋体，宽 241pt 居中并加下划线。
%   日期 小四宋体居中
%
% 关于文字：附件 3 封面给的是占位说明，此处**逐字照录**，以保证与附件 3 一致。
% 实际使用时把说明文字删掉换成自己的题目；题目字数多时自然会折成两行
% （附件 3 也注明“可分两行”）。
%
% 日期：附件 3 原文为“年      月（小四宋体，居中）”，“年”“月”之间留空隙。
% TeX 会把连续空格压成一个，故文档类提供 \gap（= \hspace{4em}）显式给出该空隙。
% 若日期是连续文字（如“二〇二六年六月”），直接写 \cdate{二〇二六年六月} 即可。
% =====================================================================

\ctitle{中文论文题目}
\etitle{English Thesis Title}
\cschool{学院名称}
\cmajor{专业班级}
\cauthor{学生姓名}
\cstudentid{学生学号}
\cadvisor{指导教师姓名}
\cadvisorinfo{指导教师职称}
\cdate{20XX年\hspace{1cm} XX月}



%% ---------- 前置部分 ----------
%% 封面 + 原创性声明；页码自原创性声明起用大写罗马数字（Ⅰ、Ⅱ、Ⅲ）
\ujsfrontmatter

% !TEX root = ../main.tex
% =====================================================================
% 中文摘要
%
% 以下文字与附件 3 中文摘要页完全一致，逐字照录（含附件 3 自带的格式说明
% 段落）。示例中“摘要：”之后为说明文字，实际使用时请替换为论文摘要正文
% （规范要求不少于 400 字）。
%
% 环境参数：#1 = 摘要页题目（黑体三号，居中，段前 1 行、段后 0.5 行）
%           #2 = 关键词（词间用中文分号“；”，末位不加标点）
% =====================================================================

\begin{cnabstract}{中文论文题目}{关键词1；关键词2；关键词3；关键词4；关键词5}

{\heiti\zihao{4}摘要}：摘要是论文的内容不加注释和评论的简短陈述。摘要应具有独立性和自含性，即不阅读论文的全文，就能获得必要的信息。摘要中有数据、有结论，是一篇完整的短文，可以独立使用。摘要应说明研究工作目的、实验方法、结果和最终结论等，重点是结果和结论。中文摘要一般不少于400字。如遇特殊需要，字数可略多。摘要中一般不用图、表、化学结构式、非公知公用的符号和术语。

摘要论文题目黑体三号，分1或2行居中打印，段前间距为1行，段后间距为0.5行；论文题目下“专业班级、学生姓名、指导教师、职称”，居中，宋体小四号，1.5倍行距。

“摘要”二字用黑体四号，1.5倍行距，段前0.5行，首行缩进2字符。“摘要”二字后打印摘要内容，内容用宋体小四号，1.5倍行距，首行缩进2字符。

关键词是为了文献标引工作从论文中选取出来的以表示全文主题内容信息款目的单词或术语。每篇论文选取3-5个词作为关键词，以显著的字符另起一行，排在摘要的左下方。尽量用《汉语主题词表》等词表提供的规范词。

关键词段前间距为1行，左顶格。“关键词”三字用黑体四号，“关键词”后为内容，用宋体小四号，每两个关键词之间用中文分号“；”间隔，最后一个关键词后不加标点符号。

从原创性声明开始设置页眉，原创性声明、中文摘要、英文摘要、目录页的页眉分两栏，左侧以小五号宋体分别键入“原创性声明”、“摘要”、“ABSTRACT”、“目录”，右侧统一键入“江苏大学本科毕业设计（论文）”。

从正文开始，以小五号宋体键入“江苏大学本科毕业设计（论文）”，靠右侧显示。原创性声明、中文摘要、英文摘要、目录页的页码使用罗马数字（Ⅰ、Ⅱ、Ⅲ）编排；正文页码从绪论开始至附录（正文-致谢-参考文献-附录）使用阿拉伯数字编排，小五号Times New Roman居中显示。

\end{cnabstract}
      % 中文摘要（含关键词）
% =====================================================================
% 英文摘要
%
% 以下文字与附件 3 英文摘要页完全一致，逐字照录（含附件 3 自带的格式说明
% 段落）。示例中“ABSTRACT”之后为说明文字，实际使用时请替换为与中文摘要
% 对应的英文摘要（规范要求不少于 350 个实词）。
%
% 环境参数：#1 = 英文题目（16pt Times New Roman，加粗，居中）
%           #2 = 英文关键词（词间用英文分号，末位不加标点）
% =====================================================================

\begin{enabstract}{英文论文题目Title}{Keywords1; Keywords2; Keywords3; Keywords4; Keywords5}

% 首段的“ABSTRACT”前缀由 UJSThesis2026.cls 的 enabstract 环境自动加（加粗、
% 首行缩进 1 字符），故此处不要再写，否则会重复成 “ABSTRACT ABSTRACT”。
英文题目、摘要和关键词内容均应与中文相同,不少于350个实词。英文题目、摘要、关键词合打一页。

所有的英文和数字均用Times New Roman字体打印。英文论文题目为四号，加粗，居中打印，段前段后间距为1行。“ABSTRACT”为小四号，加粗，首行缩进1字符打印。后空一格打印英文摘要内容，内容为小四号，1.5倍行距。

“KEY WORDS”为小四号，加粗，左起打印。后空一格打印英文关键词内容，内容为小四号，1.5倍行距，段前间距为1行。关键词之间用英文“;”间隔，最后一个关键词后不加标点符号。

\end{enabstract}
    % 英文摘要（含关键词）

%% 目录（列至三级）
\ujscontents

%% 可选：插图和附表清单
% % =====================================================================
% 插图和附表清单（可选，见规范 一·7）
%   位于目录之后；由 \listoffigures / \listoftables 自动生成。
%   同时列入目录；页码续用前置部分的大写罗马数字。
%
%   注意：LaTeX 的清单标题默认取 \listfigurename / \listtablename，
%   与“插图清单 / 附表清单”不一致，故在下面重新定义。
% =====================================================================

\clearpage
\ujsheadleft{插图和附表清单}%
\thispagestyle{ujs@front}%
\phantomsection
\addcontentsline{toc}{chapter}{插图和附表清单}

\begingroup
\renewcommand{\listfigurename}{插\hspace{0.5em}图\hspace{0.5em}清\hspace{0.5em}单}
\renewcommand{\listtablename}{附\hspace{0.5em}表\hspace{0.5em}清\hspace{0.5em}单}
\renewcommand{\cftloftitlefont}{\hfill\heiti\zihao{2}}
\renewcommand{\cftlottitlefont}{\hfill\heiti\zihao{2}}
\renewcommand{\cftafterloftitle}{\hfill}
\renewcommand{\cftafterlottitle}{\hfill}
\setlength{\cftbeforeloftitleskip}{1\baselineskip}
\setlength{\cftbeforelottitleskip}{1\baselineskip}
\setlength{\cftafterloftitleskip}{0.5\baselineskip}
\setlength{\cftafterlottitleskip}{0.5\baselineskip}
\listoffigures
\vspace{2\baselineskip}
\listoftables
\endgroup
\clearpage

%% 可选：注释表（符号、标志、缩略词、单位、术语、名词等）
% % =====================================================================
% 注释表（可选，见规范 一·8）
%   含符号、标志、缩略词、首字母缩写、单位、术语、名词等。
%   正文小四宋体，1.5 倍行距。
% =====================================================================

\clearpage
\ujsheadleft{注释表}%
\thispagestyle{ujs@front}%
\phantomsection
\addcontentsline{toc}{chapter}{注释表}

\begin{center}
  \vspace{0.5\baselineskip}
  {\heiti\zihao{-2}注\hspace{1em}释\hspace{1em}表}
\end{center}
\par\vspace{\baselineskip}

\noindent
{\songti\zihao{-4}
  \renewcommand{\arraystretch}{1.3}%
  \begin{tabularx}{\textwidth}{@{}lL@{}}
    \toprule
    \textbf{符号 / 缩略词} & \textbf{含义} \\
    \midrule
    UJS      & Jiangsu University，江苏大学 \\
    GB/T     & 中华人民共和国推荐性国家标准 \\
    LaTeX    & 一种基于 \TeX{} 的排版系统 \\
    xeCJK    & 在 XeLaTeX 下处理中日韩文字的宏包 \\
    ctex     & 面向中文排版的 LaTeX 宏包与文档类集合 \\
    bst      & BibTeX 样式文件（bibliography style） \\
    \bottomrule
  \end{tabularx}
}
\clearpage


%% ---------- 正文（从引言起用阿拉伯数字页码） ----------
\ujsmainmatter
% !TEX root = ../main.tex
% 引言示例，正式写作时请替换为论文实际内容。
\ujsunnumberedchapter{引\qquad 言}[引言]

本文围绕所选课题开展研究，介绍研究背景、研究目的与范围，概述相关领域的研究基础及尚待解决的问题。此处为引言占位文字，请按实际课题替换。

引言应说明研究目的与范围，介绍相关领域的研究基础、尚待解决的问题以及本文的组织安排。此处为模板示例文字，正式写作时应结合实际文献和研究内容补充。

  % 引言为第 1 页，后续连续编号

% !TEX root = ../main.tex
% =====================================================================
% 正文第 1 章
%
% 以下文字与附件 3 正文第 1 章完全一致，逐字照录（含附件 3 自带的格式说明
% 段落）。章标题、节标题、小节标题的分级依附件 3 所标示的字号：
%   黑体小二号（18pt）→ 章标题
%   宋体四号（14pt）  → 节标题（\section）
%   宋体小四（12pt）  → 小节标题（\subsection）
%
% 章标题为附件 3 原文“1  **”，故用不编号章命令原样输出，
% 以保证目录条目与正文标题和附件 3 逐字一致。
% =====================================================================

\ujschapter{1\quad 绪论}

\section{课题背景和意义}

\subsection{研究背景}

说明本课题的意义、目的、研究范围及要达到的技术要求；简述本课题在国内外的发展概况及存在的问题；说明本课题的指导思想；阐述本课题应解决的主要问题。

\subsection{研究意义}

……………………（正文小四，中文宋体，英文数字Times New Roman，1.5行距，首行缩进2字符）

\section{课题研究现状}

\section{本文研究内容}

正文部分：章标题为黑体小二号字体，居中打印，段前、段后间距为1行。

二级标题为宋体四号字，加粗，左对齐，段前、段后间距为0.5行。

三级标题为宋体小四号字，加粗，左对齐，段前、段后间距为0.5行。

正文中汉字用宋体小四号，英文和数字用Times New Roman小四号，1.5倍行距，首行缩进2字符。每一个新的章节均另起一页开始书写。

% !TEX root = ../main.tex
% 格式测试用引用，文献来自官方模板示例，与论文主题无关。
\section{参考文献引用格式测试}
以下段落仅用于检查 BibTeX 引用和参考文献排版，正式写作时可删除本节及对应的 input。

单篇引用测试：中文专著的引用应显示为上标方括号编号\upcite{chen2000}。再次引用同一文献\upcite{chen2000}，编号应保持一致。

连续多篇引用测试：一次引用三篇文献\upcite{harrison2012,yu2019,zhongguo1985}，检查连续编号是否合并为范围。

非连续引用测试：引用已经出现过的第一篇与第三篇文献\upcite{yu2019,chen2000}，检查编号是否按升序排列，并用逗号分隔。

平排引用测试：讨论文献\cite{chen2000}时，编号应与正文处于同一基线；句末上标引用应位于标点之前\upcite{chen2000}。

文献类型测试：英文会议论文集\upcite{babu2014}，中文报纸\upcite{lidalu1998}，英文报纸\upcite{french1987}，中文期刊\upcite{chen1996}，英文期刊\upcite{caplan1993}。

学位论文和报告测试：中文学位论文\upcite{zhang2020}，英文学位论文\upcite{calms1965}，中文报告\upcite{feng1997}，机构作者报告\upcite{who1970}，在线报告\upcite{guowuyuan2013}。

其他类型测试：标准\upcite{quanguo2010}，专利\upcite{liujialin1993}，专著析出文献\upcite{cheng1999}，会议析出文献\upcite{jia2011}，电子公告\upcite{guowuyuan2014,xiaoniu2001,dublincore2012}。参考文献表应按首次引用顺序列出全部 21 条文献，并检查长标题和 URL 的换行。

      % 1  **（附件 3 第 1 章）
% =====================================================================
% 正文第 2 章
%
% 附件 3 只在目录中给出该章条目“2  黑体四号”，正文部分未给出示意文字，
% 故章标题照录为“2  **”。实际使用时按同样格式自行增补章节内容。
% =====================================================================

\ujschapter{2  **}

% !TEX root = ../main.tex
% 图片格式测试，正式写作时可删除本文件的 input。
\section{图片编号与图注测试}

蓝色示例见图~\ref{fig:test-blue}，橙色示例见图~\ref{fig:test-orange}。图号按全文连续排列，图号与简短题名一并置于图片下方。

\begin{figure}[htbp]
  \centering
  \includegraphics[width=0.72\textwidth]{test-blue.pdf}
  \caption{蓝色纯色图片示例}
  \label{fig:test-blue}
\end{figure}

图~\ref{fig:test-blue}为矢量 PDF，缩放后边缘应保持清晰。再次引用同一图片，编号应保持不变。

\begin{figure}[htbp]
  \centering
  \includegraphics[width=0.60\textwidth]{test-orange.pdf}
  \caption{橙色纯色图片示例}
  \label{fig:test-orange}
\end{figure}

\section{正文图片引用测试}

图~\ref{fig:test-blue}和图~\ref{fig:test-orange}分别显示两种纯色版块。正文引用使用行内图号，不使用参考文献上标形式。第三章仍使用这些标签引用，检查跨章引用是否一致。

      % 2  **
% =====================================================================
% 正文第 3 章
%
% 附件 3 只在目录中给出该章条目“3  黑体四号”，正文部分未给出示意文字，
% 故章标题照录为“3  **”。实际使用时按同样格式自行增补章节内容。
% =====================================================================

\ujschapter{3  **}

% !TEX root = ../main.tex
% 数值为虚构的排版测试数据，正式写作时可删除本文件的 input。
\section{数值表格测试}

表~\ref{tab:test-numeric}给出三组模拟数值，表号与题名位于表上。同列使用统一小数位数并右对齐，检查整数位和小数点的纵向对齐。

\begin{table}[htbp]
  \centering
  \caption{模拟测量结果}
  \label{tab:test-numeric}
  {\songti\zihao{-4}
  \renewcommand{\arraystretch}{1.2}
  \setlength{\tabcolsep}{12pt}
  \begin{tabular}{lrrr}
    \toprule
    组别 & 电压 / V & 电流 / A & 时间 / s \\
    \midrule
    方案 A & 12.50 & 0.80 & 2.10 \\
    方案 B &  9.75 & 1.25 & 3.40 \\
    方案 C & 10.00 & 0.95 & 1.85 \\
    \bottomrule
  \end{tabular}\par}
  \vspace{3pt}
  \begin{minipage}{0.80\textwidth}
    \songti\zihao{5}\setstretch{1.0}
    注：数值仅用于排版测试；V、A、s 分别表示伏特、安培和秒。
  \end{minipage}
\end{table}

表~\ref{tab:test-numeric}各单元格均填写具体数值，未使用“同上”“同左”等替代词。再次引用该表时，编号仍为表~\ref{tab:test-numeric}。

\section{文字表格与跨章引用测试}

表~\ref{tab:test-layout}记录检查项目。图片引用仍指向第二章的图~\ref{fig:test-blue}和图~\ref{fig:test-orange}，用于验证跨章引用。

\begin{table}[htbp]
  \centering
  \caption{图表格式检查项目}
  \label{tab:test-layout}
  {\songti\zihao{-4}
  \renewcommand{\arraystretch}{1.2}
  \begin{tabularx}{0.95\textwidth}{@{}lL@{}}
    \toprule
    检查项目 & 预期显示 \\
    \midrule
    图号与图题 & 图号全文连续，简短题名置于图片下方。 \\
    表号与表题 & 表号全文连续，简短题名置于表格上方。 \\
    数值对齐 & 同列小数位统一，数值上下对齐。 \\
    正文引用 & 使用行内图号、表号，重复引用保持编号一致。 \\
    \bottomrule
  \end{tabularx}\par}
\end{table}

两张表依次编号为表~\ref{tab:test-numeric}、表~\ref{tab:test-layout}，不带章号前缀；图与表分别计数。

      % 3  **
% =====================================================================
% 正文第 4 章
%
% 附件 3 只在目录中给出该章条目“4  黑体四号”，正文部分未给出示意文字，
% 故章标题照录为“4  **”。实际使用时按同样格式自行增补章节内容。
% =====================================================================

\ujschapter{4  **}

% !TEX root = ../main.tex
% 公式排版测试；正式写作可删除对应 input。
\section{单行公式与编号引用}

单行公式见式\eqref{eq:test-ohm}。公式在全文中连续编号，编号位于公式所在行最右侧。
\begin{equation}
  U = I R
  \label{eq:test-ohm}
\end{equation}
其中，$U$、$I$、$R$ 分别表示电压、电流和电阻。再次引用式\eqref{eq:test-ohm}时，编号应保持不变。

\section{多行公式与末行编号}

式\eqref{eq:test-polynomial}为续行公式，使用 \texttt{multline} 环境，将唯一编号放在末行右侧。
\begin{multline}
  p(x)=a_0+a_1x+a_2x^2+a_3x^3+a_4x^4 \\
       {}+a_5x^5+a_6x^6+a_7x^7+a_8x^8.
  \label{eq:test-polynomial}
\end{multline}

对齐推导见式\eqref{eq:test-derivation}，前一行不编号，最后一行编号。
\begin{align}
  (a+b)^2 &= (a+b)(a+b) \notag\\
          &= a^2+2ab+b^2.
  \label{eq:test-derivation}
\end{align}

\begin{equation*}
  \frac{\mathrm{d}}{\mathrm{d}x}x^2 = 2x
\end{equation*}
以上未编号公式不消耗序号；下一条公式见式\eqref{eq:test-mean}。
\begin{equation}
  \bar{x}=\frac{1}{n}\sum_{i=1}^{n}x_i
  \label{eq:test-mean}
\end{equation}

交叉引用汇总：式\eqref{eq:test-ohm}、式\eqref{eq:test-polynomial}、式\eqref{eq:test-derivation}、式\eqref{eq:test-mean}。正文公式采用行内括号引用；附录公式见式\eqref{eq:appendix-a1}和式\eqref{eq:appendix-b1}。

      % 4  **
% !TEX root = ../main.tex
% =====================================================================
% 结论（附件 3 目录中为“5  结论（黑体四号）”）
%
% 以下文字与附件 3 第 5 章完全一致，逐字照录（含附件 3 自带的格式说明
% 段落）。章标题照录为“5  结论”。
% =====================================================================

\ujschapter{5\quad 结论}

\section{***（宋体四号，加粗，1.5倍行距，段前段后0.5行）}

结论应集中反映作者的研究成果，要求语言精炼、表述准确且完整，着重阐述自己的创造性成果及其在本研究领域中的意义、作用，同时应包括所得结果与已有结果的比较和本课题尚存在的问题，以及进一步开展研究的见解与建议。

\subsection{***（宋体小四，加粗，1.5倍行距，段前段后0.5行）}

……………………（正文小四，中文宋体，英文数字Times New Roman，1.5行距，首行缩进2字符）

\section{***（宋体四号，加粗，1.5倍行距，段前段后0.5行）}

\subsection{***（宋体小四，加粗，1.5倍行距，段前段后0.5行）}

……………………（正文小四，中文宋体，英文数字Times New Roman，1.5行距，首行缩进2字符）
    % 5  结论

%% ---------- 后置部分 ----------
% =====================================================================
% 参考文献
%
% 用 BibTeX + natbib 生成，样式为 gbt7714-numerical.bst（GB/T 7714-2015，
% 连续编号、方括号、上标引用）。条目写在根目录的 refs.bib 里。
%
% 编译顺序（BibTeX 不能省）：
%     xelatex main  →  bibtex main  →  xelatex main  →  xelatex main
% 在 VS Code / LaTeX Workshop 下选 “xe->bib->xe->xe” 配方即可。
%
% 正文里引用：\upcite{键名}（上标方括号，符合 GB/T 7714 与附件 3 示例），
% 需要平排方括号时用 \cite{键名}。
%
% 本页排版说明
%   * 标题“参考文献”由本页的 \bibsection（重定义为文档类的章标题格式）排出，
%     同时手工写一条目录项——因为这里不走 \chapter，latex 不会自动加目录。
%   * 条号顶格、续行缩进 2 字符（文档类已显式设置数字制列表缩进），
%     与附件 3 的版式一致。
% =====================================================================

\clearpage
\phantomsection
\addcontentsline{toc}{chapter}{参考文献}%
\begingroup
\setstretch{1.5}%
\songti\zihao{-4}%
% 对齐附件 3：取消文献条目间额外间距，正文基线约 23.3pt。
\setlength{\bibsep}{0pt}%
\setlength{\baselineskip}{23.3pt}%
\makeatletter
\let\@bibsetup\NAT@bibsetnum
\renewcommand{\bibsection}{%
  \vspace*{-3.386pt}%
  \ujs@chapterfmt{\bibname}%
  \vspace{-5.815pt}%
}
\makeatother
\bibliography{refs}
\endgroup





    % 参考文献（由 refs.bib 自动生成，列入目录）
% =====================================================================
% 致谢（可选）
%
% 以下文字与附件 3 致谢页完全一致，逐字照录（含附件 3 自带的格式说明
% 段落）。章标题照录为“致    谢”（中间空 4 格）。
%   内容中文用宋体小四号，1.5 倍行距，首行缩进 2 字符。
% =====================================================================

% 正文标题“致    谢”（中间空 4 格）；目录条目用不带空隙的“致谢”。
\ujsunnumberedchapter{致\qquad 谢}[致谢]

致谢（可选）应以简短的文字对在课题研究与设计（论文）撰写过程中直接给予帮助的指导教师、答疑教师和其他人员或单位表示自己的谢意，以示对他人劳动的尊重。语言应实事求是，切记浮夸之词。在结论后单置一页。

“致谢”二字用黑体小二号，居中，中间空4格，段前间距为2行，段后间距为1行。

内容中文用宋体小四号，1.5倍行距，首行缩进2字符。
        % 致谢（可选）
% =====================================================================
% 附录（可选）
%
% 以下文字与附件 3 附录页完全一致，逐字照录（含附件 3 自带的格式说明
% 段落）。章标题照录为“附    录”（中间空 4 格）。
%   附录的序号编排按附录 A、附录 B……编排，
%   附录内顺序可按 B2、B2.1、B2.1.1、B2.1.2 的规律编排，
%   图表按图 B1、图 B2、表 B1、表 B2 的规律编排，与正文连续编页码。
% =====================================================================

% 正文标题“附\qquad 录”（中间空 4 格）；目录条目用不带空隙的“附录”。
\ujsunnumberedchapter{附\qquad 录}[附录]

附录（可选）主要包括一些不宜放在正文中的支撑材料，如某些原始数据、数学推导、计算程序、框图、结构图、注释、统计表、计算机打印输出件等。

“附录”二字用黑体小二号，居中，中间空4格，段前间距为2行，段后间距为1行。内容中文用宋体小四号，1.5倍行距，首行缩进2字符。

附录的序号编排按附录A、附录B．．．编排，附录（以附录B为例）内的顺序可按B2，B2.1,B2.1.1,B2.1.2的规律编排。图表按：图B1，图B2，表B1，表B2的规律编排。附录与正文装订在一起，连续编页码。

% !TEX root = ../main.tex
% 附录图表公式排版测试，保留原附录说明。
\ujsappendix
\ujschapter{附录 A\quad 图表与公式测试}

本附录检查图~\ref{fig:appendix-a1}、表~\ref{tab:appendix-a1}和式\eqref{eq:appendix-a1}。编号应从 A1 开始，与正文的图~\ref{fig:test-blue}、表~\ref{tab:test-numeric}及式\eqref{eq:test-ohm}区分。

\begin{figure}[htbp]
  \centering
  \includegraphics[width=0.38\textwidth]{test-blue.pdf}
  \caption{附录蓝色图片示例}
  \label{fig:appendix-a1}
\end{figure}

\begin{table}[htbp]
  \centering
  \caption{附录 A 模拟参数}
  \label{tab:appendix-a1}
  {\songti\zihao{-4}
  \begin{tabular}{lrr}
    \toprule
    参数组 & 电阻 / $\Omega$ & 电压 / V \\
    \midrule
    A & 10.00 & 12.50 \\
    B & 15.00 &  9.75 \\
    \bottomrule
  \end{tabular}\par}
\end{table}

附录单行公式如下：
\begin{equation}
  P = U I
  \label{eq:appendix-a1}
\end{equation}
同一附录的下一条多行公式应连续编号：
\begin{align}
  P &= I(IR) \notag\\
    &= I^2R.
  \label{eq:appendix-a2}
\end{align}
再次引用图~\ref{fig:appendix-a1}、表~\ref{tab:appendix-a1}、式\eqref{eq:appendix-a1}和式\eqref{eq:appendix-a2}，编号应保持一致。

\ujschapter{附录 B\quad 编号重置测试}

进入附录 B 后，图~\ref{fig:appendix-b1}、表~\ref{tab:appendix-b1}及式\eqref{eq:appendix-b1}均应从 B1 开始。上一附录的式\eqref{eq:appendix-a2}和正文的式\eqref{eq:test-mean}仍应跳转到原位置。

\begin{figure}[htbp]
  \centering
  \includegraphics[width=0.38\textwidth]{test-orange.pdf}
  \caption{附录橙色图片示例}
  \label{fig:appendix-b1}
\end{figure}

\begin{table}[htbp]
  \centering
  \caption{附录 B 模拟参数}
  \label{tab:appendix-b1}
  {\songti\zihao{-4}
  \begin{tabular}{lrr}
    \toprule
    参数组 & 时间 / s & 电流 / A \\
    \midrule
    C & 2.10 & 0.80 \\
    D & 3.40 & 1.25 \\
    \bottomrule
  \end{tabular}\par}
\end{table}

\begin{equation}
  E = P t
  \label{eq:appendix-b1}
\end{equation}

引用汇总：图~\ref{fig:appendix-b1}、表~\ref{tab:appendix-b1}和式\eqref{eq:appendix-b1}。附录页码与正文连续，不重新从 1 开始。

      % 附录（可选）

\end{document}



